% frontmatter/zdania/verano.tex -- la página para el adulto del cuaderno
% de verano de «Zdania» (ediciones/zdania-verano.tex, ver
% notes/07-ediciones.md): la de frontmatter/verano.tex, en polaco. Va en
% vez de frontmatter/zdania/como-usar.tex, que habla del año entero.
% Detrás, en vez del mapa del año, el mapa del verano (\mapaEdicion,
% preamble.tex).
\section*{Jak korzystać z tego zeszytu na lato}

To lato z zeszytu \emph{Zdania}, drugiego z serii \emph{Uczę się czytać}
po polsku: jego ostatnie 13 tygodni, w osobnym zeszycie, żeby czytać po
trochu każdego dnia wakacji. Jest dla dziecka, które czyta już całe
zdania i chce czytać coraz płynniej -- nie uczy czytania od zera. Lato
to właśnie czas, kiedy najbardziej się przydaje, żeby nie wypaść z
rytmu.

Ma 65 stron, po jednej na każdy dzień od poniedziałku do piątku, i
każda wygląda tak samo:

\begin{itemize}[leftmargin=6mm]
\item \textbf{Data}, \textbf{Czytanie} i \textbf{Uwagi}: dla dorosłego.
  Zaznaczenie codziennie, czy dziecko czytało samodzielnie, z pomocą czy
  z trudem, nic nie kosztuje, a pod koniec lata jest zapisem jego
  postępów.
\item \textbf{Cztery zdania dnia}, w zielonej ramce: najpierw czyta się
  je po cichu, a potem na głos. Czasem są w nich dialogi, a w niektóre
  piątki (\textbf{\lblRelee}) jest tekst z całego tygodnia, żeby
  przeczytać go jeszcze raz, od początku do końca.
\item Codziennie inne \textbf{zadanie}, zawsze związane z tym, co się
  właśnie przeczytało:
  \begin{itemize}[leftmargin=6mm]
  \item \textbf{\lblResponde}: krótkie pytanie o to, co się przeczytało,
    żeby sprawdzić, że zostało zrozumiane, a nie tylko przeczytane.
  \item \textbf{\lblBusca}: znaleźć coś w tekście.
  \item \textbf{\lblRelaciona}: połączyć linią to, co do siebie pasuje.
  \item \textbf{\lblAdivina}: zagadka o historii z tego tygodnia.
  \item \textbf{\lblDibuja} i \textbf{\lblCompleta}: narysować coś, o
    czym mówi tekst, albo zamienić kilka linii w rysunek -- nie ma
    „dobrego” rysunku.
  \item \textbf{\lblTraza}: poprowadzić palcem albo ołówkiem po kropkach
    litery, wielkiej i małej, a potem napisać ją samodzielnie na linii.
  \item \textbf{\lblRepasa}: co dwa tygodnie, krótkie podsumowanie z
    kratkami do zaznaczenia i zadaniem na koniec.
  \end{itemize}
\end{itemize}

Strony opowiadają lato \textbf{Lucíi}, jej brata \textbf{Daniego} i ich
psa \textbf{Toby'ego} na wsi u babci \textbf{Rosy}: rzeka, ogród, burza,
plaża, wiejski festyn, nowy przyjaciel -- a na koniec powrót do miasta i
do szkoły. Każdy tydzień to nowy rozdział tej samej historii, tej samej
co w hiszpańskim zeszycie na lato; na następnej stronie są wszystkie, ze
słońcem do pokolorowania za każdy dzień czytania.

Wystarczy pięć albo dziesięć minut dziennie, i nie trzeba kończyć
zadania tego samego dnia -- najważniejsze jest czytanie. Jeśli dziecko
zacina się prawie na każdym słowie, niech przeczyta tylko pierwsze
zdanie, a resztę przeczyta na głos dorosły: najważniejsze jest, żeby
czytanie było przyjemnością, a nie walką. Co dziesięć stron jest mały
napis, żeby to uczcić, a na końcu zeszytu są \textbf{odpowiedzi} do
zadań \lblAdivina\ i \lblRelaciona\ oraz dyplom.

{\small\color{colorGris} Zeszyt \emph{Zdania} jest też na cały rok, z
260 stronami, a przed nim jest \emph{Pierwsze słowa}, pierwszy zeszyt z
serii.\par}
\newpage
